\subsection*{Separability, dependency audit, and further levels}

\statement{B.35}{Proved}{A sufficient product-separability condition with its quantifiers}
Let $G$ be residually finite and $H,J\leq G$. Suppose that for
every finite-index normal subgroup $U\trianglelefteq G$ the product
$(H\cap U)(J\cap U)$ is closed in the profinite topology of $G$.
Then
\[
             \overline H\cap\overline J=\overline{H\cap J}
                         \quad\hbox{inside }\widehat G.
\]
Applied to $H_T,H_5$ in B.28, this condition would prove
$\mathcal K=1$. B.52 proves the unrefined product $H_TH_5$ is closed.
The refined condition above remains unresolved; B.46--B.49
and B.52 explain exactly why the unrefined result does not
yet establish it.
\begin{proof}
Let $a\in\overline H\cap\overline J$, and fix $U$.
Choose $h\in H,j\in J$ whose images modulo $U$ agree with that
of $a$. Density in the corresponding open cosets gives
\[
 h^{-1}a\in\overline{H\cap U},\qquad
 a^{-1}j\in\overline{J\cap U}.
\]
Products of compact sets are compact, and the product of two
dense subsets is dense in their product. Hence
$h^{-1}j$ lies in the closure of $(H\cap U)(J\cap U)$.
By the hypothesis it equals $uv$ for some
$u\in H\cap U,v\in J\cap U$.
The element $hu=jv^{-1}$ belongs to $H\cap J$ and agrees
with $a$ modulo $U$. This holds for every $U$, proving the
nontrivial inclusion; the other inclusion follows from continuity.
The proof explicitly requires the refined products. It does not
replace them by the one unrefined product without another argument.
No profinite tree or complex with the required unique $A_5$ fixed
point has been constructed here. The discrete fixed Teichm\"uller
point alone supplies no such profinite object.
\end{proof}

\statement{B.36}{Proved}{The audited centralizer dependencies do not use the nonsolvable equality}
The finite-subgroup comparisons used in LM6 C.16--C.17 are
involution comparisons. The torsion-centralizer comparisons in
LM7\_PUSH E.64 are cyclic-factor comparisons; its sewing comparison
uses the specified twist stabilizer and genus-two congruence
identification. ANABELIAN\_149\_PUSH A.29l uses the separate
orbifold theorem of Boggi--Zalesskii (2019), Lemma 3.1.
None of these passages uses the unproved nonsolvable
closure-intersection equality in B.1.
\begin{proof}
For LM6 C.16--C.17 the named finite action is the single
real involution; the independently reconstructed diagrams and
centralizer argument are E.4--E.6 below. The finite comparison
is Theorem A at $p=2$, whose proof first treats prime order
by the high-degree cohomology comparison and then inducts
through the centers of finite $p$-groups. This is a comparison
for one finite $p$-subgroup, not an intersection over different
prime-power centralizers.

In LM7\_PUSH E.64.7--E.64.8 the precise groups needed are
$C_{\widehat B_3}(s)=\overline{\langle s\rangle}$ and
$C_{C_4\amalg C_2}(\bar d)=\langle\bar d\rangle$.
The proof explicitly cites Herfort--Ribes, Theorem B', and
then checks intersection with the free Ferrari kernel.
The sewing formula E.64.3 separately uses the oriented twist
stabilizer exact sequence in Boggi's automorphism paper,
Theorem 2.3 and Proposition 2.6(i). A multitwist stabilizer
is a different input from the disputed finite-subgroup
comparison. Thus the phrase ``only the element case'' should
be read as a statement about the finite centralizer dependency,
not as a claim that every ingredient of E.64 is a torsion theorem.

We retrieved the actual ANABELIAN\_149\_PUSH paper and read
A.29l on pages 115--117; it was not among the supplied attachments.
Its geometric extension is the pro-$\mathscr S$ completion of a
hyperbolic Fuchsian group, where $\mathscr S$ is the class of
finite groups of order prime to the residue characteristic.
The exhibited torsion-free surface kernel has quotient in
$\mathscr S$. These are exactly the hypotheses of
\citet[Lemma 3.1]{BoggiZalesskii2019}.
The lemma gives malnormality and self-normalization of the
finite cyclic decomposition groups and places every nontrivial
finite subgroup in one of them. Its proof uses the
$\mathscr S$-solenoid: mod-$\ell$ acyclicity forces a finite
$\ell$-subgroup to fix a unique orbifold point. It then treats
solvable finite groups by induction; for an arbitrary finite
group all its Sylow subgroups are cyclic, hence the finite
group is solvable. This proof is independent of the later
mapping-class closure-intersection assertion.
In A.29l, the geometric centralizer of a point section first
injects into the finite deck group by the separate Frobenius
weight argument. Lemma 3.1 places that finite centralizer in a
unique inertia group; the ensuing arithmetic normalizer is
proved directly by matching the arithmetic image. No centralizer
of an arbitrary nonsolvable mapping-class subgroup is substituted.
This verifies the requested dependency; it does not reassert
all the other arithmetic consequences of that return without
their own hypotheses.

The attached LM4 return already restricts F.1--F.3 to the
solvable range where appropriate and labels F.7 Open.
Its subsequent F.8--F.33 use named cyclic, dihedral, tetrahedral,
or involution comparisons, configuration maps, and finite-image
calculations, while retaining the residual $A_5$ equality as
unresolved. LM4 D.31 uses only the prime-power topology input.
LM6 C.13--C.19 use real involutions. The attached LM7\_PUSH
sections and the LM7 tower sections contain no application of
the nonsolvable mapping-class equality.
Accordingly the list of attached \emph{proved claims whose
proof requires that equality for a nonsolvable subgroup is empty}.
The places quoting it as an unresolved assertion, and the two
published sources which assert it, are recorded separately in
\texttt{records/centralizer\_dependency\_audit.md}.
This is a dependency audit of the supplied versions, not a
claim that every unrelated statement in those versions is correct.
\end{proof}

\statement{B.37}{Proved}{The level-six control also has trivial residual kernel}
For $N=2,3,4,6$, the group $\mathcal K_N$ of task B.5 is trivial.
For $N=6$ this uses the full profinite mapping class group of the
genus-one curve $Y(6)$ with its twelve cusps; it is not a
prime-by-prime linear computation.
\begin{proof}
The levels three and four were proved in B.25, and level two
is the defining three-cusp comparison: a pure multiplier-one
geometric outer action on a three-puncture sphere is trivial.
For level six, the Chinese remainder map gives
\[
 Q_6=\mathrm{PSL}_2(\mathbb Z/6)
           \simeq S_3\times A_4,
 \qquad [B:\Gamma(6)/\{\pm1\}]=72.
\]
The kernel is the intersection of the level-two and level-three
kernels. The level-two kernel is characteristic under the
peripheral group, and the level-three kernel is characteristic
under the same group: an epimorphism from $C_2*C_3$ to $A_4$
sends the two factors to orders two and three, their product
has order three, and the triangle presentation $(2,3,3)$
has order twelve. Thus all such epimorphisms have the same kernel.
An element $\Phi\in\mathcal H$ acts trivially on the outer
quotient of $S_3$. Its action on $A_4$ is inner as well:
it preserves the conjugacy class of the named three-cycle
cusp because its multiplier is one, whereas the nonidentity
outer automorphism of $A_4$ exchanges the two three-cycle
classes. After an inner adjustment in $\widehat B$, its
induced action on $Q_6$ is therefore the identity.

The level-six surface group $L$ is free of rank thirteen.
The cusp order in $Q_6$ is six, so it has twelve cusps;
$\chi(L)=72(-1/6)=-12$ gives compactified genus one.
Restriction of the adjusted automorphism to $\widehat L$
centralizes the deck action of $Q_6$ in $\Out(\widehat L)$.
If $\Phi\in\mathcal K_6$, that outer action lies in the
image of the completed genus-one mapping class group.
Asada's full genus-one congruence theorem\citep[Theorem 3A]{Asada2001} identifies that
image with the full completion. Consequently it is an
actual element of $C_{\widehat\Gamma_{1,[12]}}(Q_6)$.

The group $Q_6$ is solvable. The noncircular prime-power
normalizer induction of B.30 works for genus at most two:
the decorated quotient genus does not increase, and the
same goodness and congruence inputs apply at each stage.
Hence its completed centralizer is the closure of its
discrete centralizer. The discrete quotient orbifold has
one point of order two, one of order three, and one cusp.
Each is distinguished, so its pure mapping class group
is trivial. The covering normalizer sequence makes the
discrete centralizer its kernel $Z(Q_6)=1$.
Thus the restricted action on $\widehat L$ is trivial.
If an automorphism of $\widehat B$ is inner on the open
centerless subgroup $\widehat L$, adjust it to be identity
there. For $b\in\widehat B$, the element $b^{-1}\Phi(b)$
centralizes $\widehat L$, and that centralizer is trivial.
Thus $\Phi=1$ as an outer automorphism, proving the claim.
The exact finite control calculations are in
\texttt{code/check\_level6.py} and its output.
\end{proof}

\statement{B.38}{Open}{The remaining congruence levels}
No determination of $\mathcal K_N$ for all $N$ is obtained here.
In particular $\mathcal K_5=1$ is the main unresolved assertion,
and the corresponding level-seven and level-eleven kernels
remain undecided. Their quotient compactifications have genera
three and twenty-six, respectively. An index-seven or
index-eleven resolvent alone supplies neither an injective
comparison back to the original group nor a simultaneous
lifting theorem inside the profinite mapping class group.
For levels with higher genus one must additionally distinguish
the image in the loop-outer group from the abstract full
mapping class completion and its congruence kernel.
Even when $N\mid M$, this paper does not assert
$\mathcal K_M\subseteq\mathcal K_N$ without proving the
required equivariant profinite covering-descent comparison.

\statement{B.39}{Open}{The nonsolvable direction and the published assertions}
For a finite solvable subgroup of a genus-zero mapping class
group, the centralizer is the closure of the discrete centralizer,
by the prime-power normalizer induction in B.30 and the
open-kernel argument B.2. This proves the solvable direction
of the refined triangle-orbifold statement in task B.6.
For $A_5$ on twelve vertices the reverse direction would
assert $\mathcal K\ne1$, which remains undecided.
Accordingly Conjecture $2^-$ is not proved or refuted here.
Neither Boggi's Proposition 2.4 and Theorem C(iii), nor
Boggi--Zalesskii's Theorem 5.1(iii), is refuted by this audit.
Their nonsolvable centralizer conclusions remain unsupported
by the displayed proofs at the step B.1 isolates.
Boggi's Theorem A in its general ramified nonsolvable case
also retains the missing input in Lemma 3.2, where the
configuration-group centralizer is identified using
Proposition 2.4. The free-action proof in Section 3.1 and
the prime-power repairs above do not use that missing step.
For a particular action whose discrete centralizer is already
trivial, its own congruence property is tautological; that
observation does not prove the general ramified theorem or
the ambient profinite centralizer comparison.

\statement{B.40}{Open}{The complete relative images and the bad-prime refinements remain undetermined}
B.20 proves the stated coprime pro-$\ell$ invariant construction
and a lift to the full peripheral group. It does not determine
the full relative images requested at $U=K$ or at
$U=K'K^p$ for $p=2,3,5$. Membership in the actual profinite
mapping class image is equivalently the normalizer or
three-coset equation in B.9. B.54 now constructs relative
data at $U=K$, at the good prime seven, for which every
full peripheral R2 lift has a nontrivial tripod action.
None is a braid lift. Thus the universal assertion that
R2 can always be carried out inside that braid image is
false. This obstruction does not classify the two complete
relative images or their intersection, nor their
compatibility over finer characteristic levels.
No computation in this paper suppresses the primes
$2,3,5$ or upgrades coprime invariant lifting to those primes.

\statement{B.41}{Open}{Parity and further-orbit lifting in the hypothetical nontrivial kernel}
If $\mathcal K\ne1$, neither the restriction of $\Delta$
to all its elements nor its extension along every required
icosahedral marked-orbit locus has been determined.
The latter extension would have to be constructed with
compatible based maps and then checked against the pentagon;
it does not follow from centralization of $A_5$ at one level.
B.10 does prove that a nonidentity element of $\mathcal K$
which also satisfied the pentagon would lie outside the
Galois image. No such element or complete lifting chain is
constructed here. This conditional implication is retained
without using it as evidence that either hypothesis holds.
